% Bilan : critères E.1, limites, difficultés, perspectives, conclusion
% (docs/retours/partie-5-lot4_2026-10-31.md ; JOURNAL_ERREURS.md).

\chapter{Bilan}
\label{chap:bilan}

\section{Critères de réussite du MVP}

Le rapport de cadrage (E.1) fixait six critères ; certains ont été reformulés par des ADR avant toute
mesure.

\begin{center}
\begin{tabularx}{\textwidth}{Xll}
\toprule
Critère & État & Preuve \\
\midrule
1. Tout se reconstruit depuis le brut par des commandes documentées & rempli & 2 chargements identiques \\
2. Le modèle bat B1, intervalle à 95~\% qui exclut 0 & rempli & validation et test scellé \\
3. Calibration de \(P(T > 2{,}5)\) à \(\pm 3\) points par décile & diagnostic & ADR-0037 \\
4. Couverture de l'intervalle à \(\pm 3\) points de l'annoncée & rempli & 0,2 point au test \\
5. Base, API et interface démarrent ; une prédiction s'affiche & rempli & rejeu et live réels \\
6. Documentation et chapitres mathématiques & rempli & ce rapport, docs/ \\
\bottomrule
\end{tabularx}
\end{center}

Le critère 3 n'est pas tenu à la lettre au test scellé (5 déciles sur 10 au-delà de 3 points), alors que
la pente vaut 1,00 : avec environ 209 matchs par décile, 3 points sont sous l'erreur d'échantillonnage.
Le protocole l'avait classé en diagnostic avant toute évaluation ; il est publié tel quel.

\section{Critères de la version intermédiaire}
\label{sec:criteres-e2}

La version intermédiaire (jalon J9) répond à la question \og que rapporte vraiment la composition des
équipes ? \fg{}. Le rapport de cadrage (E.2) fixait ses critères.

\begin{center}
\begin{tabularx}{\textwidth}{Xll}
\toprule
Critère & État & Preuve \\
\midrule
Chaque groupe de variables a un apport mesuré, gain et intervalle, publié & rempli & tableau des
  ablations, chap.~\ref{chap:resultats} \\
Décision argumentée sur la stabilité du onze & rempli & abandonnée pour H2 \\
Régularisation et boosting essayés & rempli & aucune ne remplace M3 \\
Second horizon avec composition, comparé au premier, avec page d'apport & rempli en rejeu & H1 et H2 côte
  à côte, page d'apport \\
Test scellé de l'horizon H2, une seule fois & rempli & non significatif, comme annoncé \\
Prédictions live tracées avant le coup d'envoi pendant au moins 8 journées & en cours (7 sur 8) & 89
  prédictions, premier relevé \\
\bottomrule
\end{tabularx}
\end{center}

Deux éléments prévus par E.2 ne sont pas faits : les services conteneurisés (Compose) et la planification
des collectes, laissés en option par l'ADR-0047. La collecte de l'API est terminée et toute tâche planifiée
passerait par une ADR. Le critère des 8 journées dépend du calendrier : il manque une journée de matchs
tracée avant le coup d'envoi.

\section{Limites}

\begin{itemize}
  \item \textbf{Un second horizon de rejeu seulement.} Avec la composition, H2 n'a pas de live : aucune
        composition n'arrive plus après la fin de l'abonnement, et l'application répond \og indisponible
        \fg{}. Il explique ce que la composition change, il ne prédit pas l'avenir.
  \item \textbf{Un gain de composition minuscule et fragile} : \(0{,}0008\) nat par match, borne basse de
        l'intervalle à \(+0{,}00003\), plus faible à fenêtre égale, concentré en début de saison. Le test
        scellé de H2 ne peut pas le distinguer de zéro (puissance de l'ordre de \(0{,}007\)). Une part de
        l'information de la composition est déjà dans l'Elo (corrélation 0,83 avec la qualité du onze).
  \item \textbf{Des variables de composition imparfaites} : notes calculées après le match par le
        fournisseur, selon une formule inconnue ; les blessures ne sont pas modélisées (un absent connu
        seulement par son absence du onze) ; une variable de G5, surtout l'âge moyen faute de date de
        naissance, manque pour 8~\% des matchs du test.
  \item \textbf{Un seul pli scellé}, d'une saison et quart, utilisé pour deux tests successifs : le résultat
        vaut pour cette période et ces cinq championnats, et l'intervalle de H2 est large (1\,923 matchs).
  \item \textbf{Approximations du live} : moyennes glissantes approchées quand le championnat n'est pas à
        jour (dans une tolérance vérifiée), tirs de football-data à la place de ceux de l'API.
  \item \textbf{Journal prospectif encore court} : 89 prédictions, 7 journées, 48 résultats sur 4 jours ;
        le premier relevé ne conclut pas, par règle.
\end{itemize}

\section{Difficultés}

Le journal compte 68 entrées. Quelques-unes ont changé la méthode :

\begin{itemize}
  \item \textbf{Le premier lancement du test scellé a échoué} (E-059) : une date de l'expérience, lue
        comme un objet date, a fait planter l'écriture du rapport. Aucun résultat n'avait été produit ;
        la relance unique a été décidée par écrit, avec ses règles de lecture, avant d'être lancée.
        Leçon : une commande qui ne s'exécute qu'une fois se répète d'abord de bout en bout, sur son
        vrai fichier d'entrée, avec des données synthétiques.
  \item \textbf{Une exactitude affirmée sans test} (E-063) : la première version de la fraîcheur par
        équipe se disait exacte ; les tests ont montré que l'a priori des moyennes glissantes dépend de
        tout le championnat. La décision a été reprise sur une mesure, avant la fusion.
  \item \textbf{Des identifiants internes instables} (E-058) : ils changent quand le brut grandit. Hors du
        référentiel, on rapproche par l'identifiant de l'API.
  \item \textbf{Un test scellé qui n'aurait pas pu s'exécuter} (E-067) : en relisant la procédure avant de
        poser le tag de gel de H2, on a vu que la construction du jeu scellé, qui vit dans le code gelé,
        était refusée jusqu'au test. Corrigé avant le tag, avec la décision écrite. Leçon : avant de poser un
        gel, parcourir \emph{toutes} les commandes de la phase qu'il protège, pas seulement l'évaluation.
  \item \textbf{Une garde qui n'en était pas une} : une fois le test passé, rien ne protégeait plus le modèle
        H2 testé, le gel étant levé. Une garde chiffrée a été ajoutée avant l'ouverture du rejeu : le rejeu
        doit reproduire les prédictions des tests (ADR-0052).
  \item \textbf{Les rapports d'expérience de 700 Ko} (E-066) : refusés par le hook des gros fichiers, ce qui
        a laissé un lancement noté \og arbre non propre \fg{} ; refait sur arbre propre, les deux conservés.
  \item \textbf{L'outillage} : intégration continue fusionnée avant la fin de ses tests (E-047), échec de
        test masqué par un tube (E-048), formatage des blocs de code des fichiers Markdown (E-049).
        Chacun a laissé une règle écrite.
\end{itemize}

\section{Perspectives}

\begin{itemize}
  \item \textbf{Huitième journée et suivantes} : poursuivre les prédictions live avant le coup d'envoi,
        refaire le relevé après chaque collecte de résultats, et n'en tirer une conclusion qu'à partir de
        8 journées.
  \item \textbf{Un live pour H2} supposerait un réabonnement ou une autre source de compositions : hors
        périmètre aujourd'hui.
  \item \textbf{Blessures} : hors du périmètre de la version intermédiaire (ADR-0047) ; la composition
        annoncée n'en voit que l'effet, par l'absence d'un joueur du onze.
  \item \textbf{Exploitation} : services conteneurisés et collecte live planifiée, laissés en option ; une
        planification de collecte passerait par sa propre ADR.
  \item \textbf{Suivi de la calibration dans le temps}, version avancée du cadrage.
\end{itemize}

\section{Conclusion}

Le MVP est atteint : une chaîne complète, reproductible depuis le brut, qui prédit une loi du nombre de
buts mieux que la référence par championnat et par terrain, sur des matchs jamais vus, avec une
incertitude annoncée qui tient ses promesses. Le gain est petit en valeur absolue (0,019 nat par match),
comme le laissait prévoir le bruit du problème, mais il est mesuré avec un intervalle et confirmé hors
échantillon. Les choix qui ont le plus compté ne sont pas des modèles : sauvegarder le brut avant la fin
de l'abonnement, figer le protocole avant de mesurer, tenir un ensemble de test sous scellés, et refuser
de prédire plutôt que d'inventer une valeur.

La version intermédiaire tient ces mêmes engagements sur une question plus étroite : que rapporte la
composition des équipes ? Réponse mesurée, avec ses intervalles : presque rien. La qualité du onze apporte
un gain de \(0{,}0008\) nat par match, de justesse, fragile et en grande partie déjà dans l'Elo ; la
stabilité du onze et l'entraîneur n'apportent rien ; ni la régularisation ni le boosting ne battent le
Poisson par équipe ; le test scellé de ce second horizon, mené une seule fois, a le même signe sans être
significatif, comme son calcul de puissance l'annonçait. Un résultat aussi modeste, publié tel quel après
des règles écrites avant de regarder, vaut mieux qu'un gain apparent obtenu en cherchant : c'est ce que le
protocole est fait pour garantir.
